\documentclass[aspectratio=169,]{beamer}

% -------------------------------------------------------------------------
% 1. REQUIRED PACKAGES (The "Batteries" Quarto usually includes)
% -------------------------------------------------------------------------
\usepackage{graphicx}
\usepackage{booktabs}
\usepackage{longtable}
\usepackage{hyperref}
\usepackage{array}      % <--- ADD THIS LINE HERE
\usepackage[most]{tcolorbox} % REQUIRED for Quarto Callouts
\usepackage{fontawesome5}    % REQUIRED for Callout Icons
\usepackage{framed}          % REQUIRED for standard blocks
\usepackage{tikz}
\usetikzlibrary{arrows.meta}

% Fix for Pandoc's "tightlist"
\providecommand{\tightlist}{%
  \setlength{\itemsep}{0pt}\setlength{\parskip}{0pt}}

% -------------------------------------------------------------------------
% 2. CUSTOM DESIGN PACKAGES
% -------------------------------------------------------------------------
\usepackage{tikz}
\usepackage{calc}
\usepackage{xcolor}
\usepackage{xparse} 
\usepackage{etoolbox} 
\usepackage[default]{sourcesanspro} 
\usetikzlibrary{calc, positioning, shapes.geometric}

% -------------------------------------------------------------------------
% DEFINE QUARTO CALLOUT COLORS (REQUIRED)
% -------------------------------------------------------------------------
% We map these to your "Elegant Dark Mode" aesthetic.
% Note: Quarto expects these EXACT names.

% 1. Note (Blue)
\definecolor{quarto-callout-note-color}{HTML}{3498DB}        % Icon/Title Color
\definecolor{quarto-callout-note-color-frame}{HTML}{2980B9}  % Border Color

% 2. Tip (Green)
\definecolor{quarto-callout-tip-color}{HTML}{27AE60}
\definecolor{quarto-callout-tip-color-frame}{HTML}{2ECC71}

% 3. Important (Red/Crimson)
\definecolor{quarto-callout-important-color}{HTML}{C0392B}
\definecolor{quarto-callout-important-color-frame}{HTML}{E74C3C}

% 4. Warning (Orange)
\definecolor{quarto-callout-warning-color}{HTML}{F39C12}
\definecolor{quarto-callout-warning-color-frame}{HTML}{E67E22}

% 5. Caution (Yellow/Gold)
\definecolor{quarto-callout-caution-color}{HTML}{F1C40F}
\definecolor{quarto-callout-caution-color-frame}{HTML}{D4AC0D}

% Add this right before \begin{document}
\newcommand{\Est}[1]{\textcolor{quarto-callout-important-color}{\mathbf{#1}}}

\providecommand{\estc}[1]{\textcolor[HTML]{1B7837}{#1}}
\makeatletter
\@ifpackageloaded{caption}{}{\usepackage{caption}}
\AtBeginDocument{%
\ifdefined\contentsname
  \renewcommand*\contentsname{Table of contents}
\else
  \newcommand\contentsname{Table of contents}
\fi
\ifdefined\listfigurename
  \renewcommand*\listfigurename{List of Figures}
\else
  \newcommand\listfigurename{List of Figures}
\fi
\ifdefined\listtablename
  \renewcommand*\listtablename{List of Tables}
\else
  \newcommand\listtablename{List of Tables}
\fi
\ifdefined\figurename
  \renewcommand*\figurename{Figure}
\else
  \newcommand\figurename{Figure}
\fi
\ifdefined\tablename
  \renewcommand*\tablename{Table}
\else
  \newcommand\tablename{Table}
\fi
}
\@ifpackageloaded{float}{}{\usepackage{float}}
\floatstyle{ruled}
\@ifundefined{c@chapter}{\newfloat{codelisting}{h}{lop}}{\newfloat{codelisting}{h}{lop}[chapter]}
\floatname{codelisting}{Listing}
\newcommand*\listoflistings{\listof{codelisting}{List of Listings}}
\makeatother
\makeatletter
\makeatother
\makeatletter
\@ifpackageloaded{caption}{}{\usepackage{caption}}
\@ifpackageloaded{subcaption}{}{\usepackage{subcaption}}
\makeatother

% -------------------------------------------------------------------------
% NEW: COLOR THEME COMMANDS (Robust Single-Slide Overrides)
% -------------------------------------------------------------------------

% 1. Light Mode Command
\newcommand{\LightMode}{
  % 1. Draw White Background (TikZ Overlay)
  \begin{tikzpicture}[remember picture, overlay]
    \fill[white] (current page.south west) rectangle (current page.north east);
  \end{tikzpicture}
  % 2. Force Text Colors to Black
  \setbeamercolor{normal text}{fg=black}
  \setbeamercolor{frametitle}{fg=black}
  \setbeamercolor{itemize item}{fg=black}
  \setbeamercolor{structure}{fg=black}
  \usebeamercolor[fg]{normal text}
}

% 2. Black Mode Command
\newcommand{\BlackMode}{
  \begin{tikzpicture}[remember picture, overlay]
    \fill[black] (current page.south west) rectangle (current page.north east);
  \end{tikzpicture}
  \setbeamercolor{normal text}{fg=white}
  \setbeamercolor{frametitle}{fg=white}
  \setbeamercolor{itemize item}{fg=white}
  \setbeamercolor{structure}{fg=white}
  \usebeamercolor[fg]{normal text}
}

% 3. Brand Mode Command
\newcommand{\BrandMode}{
  \begin{tikzpicture}[remember picture, overlay]
    \definecolor{tempBrand}{HTML}{7D181E}
    \fill[tempBrand] (current page.south west) rectangle (current page.north east);
  \end{tikzpicture}
  \setbeamercolor{normal text}{fg=white}
  \setbeamercolor{frametitle}{fg=white}
  \setbeamercolor{itemize item}{fg=white}
  \setbeamercolor{structure}{fg=white}
  \usebeamercolor[fg]{normal text}
}

% -------------------------------------------------------------------------
% 4. UC BERKELEY THEME (Modern Matte) - FIXED SCOPE
% -------------------------------------------------------------------------

% 1. DEFINE COLORS GLOBALLY (Required for assumptionbox and modes)
\definecolor{BerkeleyBlue}{HTML}{003262}
\definecolor{CaliforniaGold}{HTML}{FDB515}
\definecolor{BerkeleyMatte}{HTML}{0D2D52} 
\definecolor{SoftWhite}{HTML}{F8F9FA}

% 4. UC Berkeley Theme (Modern Matte)
\newcommand{\BerkeleyMode}{
  % 2. Draw Background AND Title
  \begin{tikzpicture}[remember picture, overlay]
    % Background Fill (Covers old title)
    \fill[BerkeleyMatte] (current page.south west) rectangle (current page.north east);
    % Gold Accent Strip
    \fill[CaliforniaGold] (current page.south west) rectangle ($(current page.north west)+(0.15cm,0)$);
    % REDRAW TITLE ON TOP
    \node[anchor=west, white, font=\Large\bfseries] at ($(current page.north west)+(1cm,-1cm)$) {\insertframetitle};
  \end{tikzpicture}
  
  % 3. Set Text Colors
  \setbeamercolor{normal text}{fg=SoftWhite}
  \setbeamercolor{frametitle}{fg=BerkeleyMatte} % Hide original title
  \setbeamercolor{itemize item}{fg=CaliforniaGold} 
  \setbeamercolor{structure}{fg=CaliforniaGold}
  
  \usebeamercolor[fg]{normal text}
}

% 5. UC Berkeley Theme (Academic Light)
\newcommand{\BerkeleyLightMode}{
  % (Color definitions removed here as they are now global)
  
  % 2. Draw Background AND Title
  \begin{tikzpicture}[remember picture, overlay]
    % Background Fill
    \fill[white] (current page.south west) rectangle (current page.north east);
    % Blue Header Bar
    \fill[BerkeleyBlue] (current page.north west) rectangle ($(current page.north east)+(0,-1.4cm)$);
    % Gold Line
    \fill[CaliforniaGold] ($(current page.north west)+(0,-1.4cm)$) rectangle ($(current page.north east)+(0,-1.45cm)$);
    % REDRAW TITLE ON TOP
    \node[anchor=west, white, font=\Large\bfseries] at ($(current page.north west)+(1cm,-0.7cm)$) {\insertframetitle};
  \end{tikzpicture}
  
  % 3. Set Text Colors
  \setbeamercolor{normal text}{fg=black}
  \setbeamercolor{frametitle}{fg=white} 
  \setbeamercolor{itemize item}{fg=BerkeleyBlue}
  \setbeamercolor{structure}{fg=BerkeleyBlue}
  
  \usebeamercolor[fg]{normal text}
}

% -------------------------------------------------------------------------
% NEW: ACADEMIC POWER LAYOUTS
% -------------------------------------------------------------------------

% 1. Statement Layout (Big Impact Transitions)
% Usage: ::: {.layout-statement} Content :::
\newenvironment{layout-statement}{
  \vspace*{\fill}
  \centering
  \Huge \bfseries
  \setbeamercolor{normal text}{fg=CaliforniaGold} % Use Gold for impact
  \usebeamercolor[fg]{normal text}
}{
  \vspace*{\fill}
}

% 2. Dense Data Layout (Regression Tables)
% Usage: ::: {.layout-dense} Table :::
\newenvironment{layout-dense}{
  \vspace*{\fill}
  \tiny                  % Force tiny font
  \setlength\tabcolsep{2pt} % Tighten table columns
  \centering
}{
  \vspace*{\fill}
}

% 3. Split Slide Macro (Vertical Centering)
% Usage: \SplitSlide{Left Content}{Right Content}
\newcommand{\SplitSlide}[2]{
  \begin{columns}[c] % 'c' option forces vertical centering
    \begin{column}{0.48\textwidth}
      #1
    \end{column}
    \begin{column}{0.48\textwidth}
      \centering
      #2
    \end{column}
  \end{columns}
}

% -------------------------------------------------------------------------
% 3. DYNAMIC BACKGROUND & CONTRAST LOGIC
% -------------------------------------------------------------------------
  \definecolor{mainbg}{HTML}{0D2D52}

\setbeamercolor{background canvas}{bg=mainbg}

% Logic to calculate luminance and set text color automatically
\makeatletter
\newcommand{\AutoSetContrast}{
  \extractcolorspec{mainbg}{\mainbgspec}
  \expandafter\convertcolorspec\mainbgspec{rgb}\mainbgrgb
  \def\calculatebrightness##1,##2,##3;{
    \pgfmathparse{0.2126*##1 + 0.7152*##2 + 0.0722*##3}
  }
  \expandafter\calculatebrightness\mainbgrgb;
  \pgfmathparse{\pgfmathresult > 0.5 ? 1 : 0}
  \ifnum\pgfmathresult=1
    \definecolor{maintext}{HTML}{222222}
    \definecolor{dimtext}{HTML}{555555}
  \else
    \definecolor{maintext}{HTML}{FFFFFF}
    \definecolor{dimtext}{HTML}{CCCCCC}
  \fi
}
\makeatother

\AutoSetContrast

\setbeamercolor{normal text}{fg=maintext}
\setbeamercolor{frametitle}{fg=maintext}
\setbeamercolor{title}{fg=maintext}
\setbeamercolor{structure}{fg=maintext} 

% -------------------------------------------------------------------------
% 4. FOOTER CUSTOMIZATION
% -------------------------------------------------------------------------
\setbeamercolor{footline}{fg=dimtext}
\setbeamerfont{footline}{size=\tiny}

\newtoggle{hidesection}
\togglefalse{hidesection}

\setbeamertemplate{footline}{
  \begin{beamercolorbox}[wd=\paperwidth, ht=2.5ex, dp=1.125ex, leftskip=.3cm, rightskip=.3cm]{footline}
    \iftoggle{hidesection}{}{
      \insertsectionhead
    }
    \hfill
    \insertframenumber
  \end{beamercolorbox}
}
\setbeamertemplate{navigation symbols}{}

% -------------------------------------------------------------------------
% 5. CUSTOM COMMANDS & ENVIRONMENTS (FIXED FOR STABILITY)
% -------------------------------------------------------------------------
\NewDocumentCommand{\navbutton}{ O{} O{fill=gray} m }{
  \hyperlink{#1}{
    \tikz[baseline=(node.base)]{
      \node(node)[
        anchor=base,
        rounded corners=3pt,
        inner sep=5pt,
        text=white,
        font=\small\bfseries,
        #2
      ]{#3};
    }
  }
}

\newenvironment{layout-math}{
  \vspace*{\fill}
  \centering
  \begin{minipage}{\linewidth}
  \centering
}{
  \end{minipage}
  \vspace*{\fill}
}

% SAFE Layout for Tables (Removed risky \makebox)
\newenvironment{layout-table}{
  \vspace*{\fill}
  \begin{center}
}{
  \end{center}
  \vspace*{\fill}
}

% SAFE Layout for Full Figures (Uses LRBOX to capture content first)
\newsavebox{\fullfigbox}
\newenvironment{layout-full-figure}{
  \begin{lrbox}{\fullfigbox}
    \begin{minipage}{\paperwidth}
      \centering
}{
    \end{minipage}
  \end{lrbox}
  \begin{tikzpicture}[remember picture, overlay]
    \node[anchor=center, inner sep=0pt] at (current page.center) {
      \usebox{\fullfigbox}
    };
  \end{tikzpicture}
}

\newcommand{\hidefootersection}{\global\toggletrue{hidesection}}
\newcommand{\showfootersection}{\global\togglefalse{hidesection}}


% -------------------------------------------------------------------------
% CUSTOM BERKELEY CALLOUT BOX
% -------------------------------------------------------------------------
\newtcolorbox{assumptionbox}[1]{
  enhanced,
  colback=SoftWhite,       % White-ish background
  colframe=CaliforniaGold, % Gold Border
  coltitle=SoftWhite,      % White Title Text
  colbacktitle=BerkeleyMatte, % Berkeley Blue Title Background
  title={#1},
  fonttitle=\bfseries\small,
  arc=2pt,                 % Slight rounded corners
  boxrule=1pt,             % Thin elegant border
  left=5pt, right=5pt, top=5pt, bottom=5pt, % Padding
  titlerule=0pt,
  width=\linewidth
}

% -------------------------------------------------------------------------
% 6. DOCUMENT BODY
% -------------------------------------------------------------------------
\title{Internalizing Environmental Risk: Insurance Design and Firm
Behavior in Hazardous Industries}
\author{Kaleb K. Javier}

\begin{document}

\frame{\titlepage}

\begin{frame}{}
\phantomsection\label{section}
\LightMode
\centering

\includegraphics[width=1\textwidth,height=\textheight]{../../Data/tank_jpg.jpg}
\end{frame}

\begin{frame}{}
\phantomsection\label{section-1}
\LightMode
\centering

\includegraphics[width=1\textwidth,height=\textheight]{../../Data/LUST.jpg}
\end{frame}

\begin{frame}
\begin{block}{Why Insurance Design Is a First-Order Policy Question}
\phantomsection\label{why-insurance-design-is-a-first-order-policy-question}
\BerkeleyLightMode

\begin{itemize}
\tightlist
\item
  The firm pays for what it causes, regardless of fault --- across oil
  \& gas (OPA 1990), surface mining (SMCRA 1977), and hazardous
  substances (CERCLA 1980, RCRA 1976)
\item
  Works when firms can be held fully accountable --- breaks down when
  they cannot (Shavell 1986, 2007)
\item
  Shavell (2004): two instruments to restore accountability:

  \begin{itemize}
  \tightlist
  \item
    Minimum Asset / bonding requirements
  \item
    \textbf{Compulsory liability insurance} --- where the premium does
    the deterrence work, \emph{if} the insurer can price on
    risk-relevant attributes
  \end{itemize}
\end{itemize}

\vfill

\centering\textbf{This paper examines the design and efficacy of compulsory liability insurance as a financial responsibility instrument in environmental hazard settings.}
\end{block}
\end{frame}

\begin{frame}
\begin{block}{Does insurance design change firm behavior?}
\phantomsection\label{does-insurance-design-change-firm-behavior}
\BerkeleyLightMode

\textbf{Setting}

\begin{itemize}
\tightlist
\item
  USTs --- retail petroleum storage, \(\sim\),140,000 facilities, the
  most common source of groundwater contamination in the US (EPA 2024)
\item
  Compulsory insurance the law since 1988 (RCRA Subtitle I, \$1M floor)
\item
  \textbf{Clean shock:} Texas exits its flat-fee fund December 22, 1998
  --- 26,000 facilities enter risk-rated private insurance overnight
\item
  9.8M tank-years of EPA panel data spanning both regimes
\end{itemize}

\vspace{0.2cm}

\textbf{Three questions}

\begin{itemize}
\tightlist
\item
  \textbf{Q1 --- Observability:} Is release risk predictable from
  observable tank characteristics?
\item
  \textbf{Q2 --- Behavior:} Does risk-based pricing change closure,
  replacement, and release outcomes?
\item
  \textbf{Q3 --- Welfare:} How much of the first-best welfare gain does
  risk-rated pricing capture?
\end{itemize}
\end{block}
\end{frame}

\begin{frame}
\begin{block}{What is an underground storage tank?}
\phantomsection\label{what-is-an-underground-storage-tank}
\BerkeleyLightMode

{[}ADD FIGURE HERE{]}
\end{block}
\end{frame}

\begin{frame}
\begin{block}{Related literature}
\phantomsection\label{related-literature}
\BerkeleyLightMode

\small

\textbf{Liability regimes, financial responsibility, and UST policy}

Shavell (1986, 2004, 2007); Yin, Kunreuther \& White (2011); Yin, Pfaff
\& Kunreuther (2011, 2007); Boomhower (2019); Ho, Hsu, Cha \& Rivera
(2018)

\emph{First firm-level panel analysis of insurance design --- risk-rated
vs.~flat-fee --- as a financial responsibility instrument; clean policy
shock; structural welfare decomposition}

\medskip

\textbf{Environmental economics and environmental risk}

Fowlie, Reguant \& Ryan (2016); Muehlenbachs (2015); Blundell,
Gowrisankaran \& Langer (2020); Kellogg \& Reguant (2021)

\emph{Dynamic structural model of firm behavior under environmental
regulation; welfare decomposition bounding gap to first-best}

\medskip

\textbf{IO of selection markets: risk-based and cost-based pricing}

Einav, Finkelstein \& Mahoney (2021); Einav, Finkelstein \& Cullen
(2010); Bundorf, Levin \& Mahoney (2012); Handel, Hendel \& Whinston
(2015); Mahoney \& Weyl (2017)

\emph{{[}PLACEHOLDER{]}}
\end{block}
\end{frame}

\end{document}
